\section{Introduction}
\label{ch:intro}

\subsection{Background}

A \emph{Systematic Investment Plan} (SIP) is the most common way Indian retail investors buy
mutual funds: a fixed amount, for example \Rs10{,}000, is invested on the same date every month.
Most SIPs go into a single equity fund. Equity has the highest long-run return of the
mainstream asset classes, but it also falls the hardest: the Nifty~50 lost about 24\% in 2011 and
about 28\% during the COVID-19 crash of March 2020 (Section~\ref{ch:results}).

Two other assets are easy for an Indian investor to hold through mutual funds and exchange-traded
funds. \emph{Gold}, bought in rupees, has historically moved largely independently of Indian equities
and often rises when the rupee falls \citep{baur2010}. A \emph{liquid fund} invests in very
short-term government and money-market paper; it earns roughly the 91-day Treasury-bill rate and
almost never falls in value. Holding all three should give a smoother ride than equity alone. The
question is \emph{how much} of each month's instalment should go into each.

\subsection{Problem statement}

\begin{quote}
\emph{Given a fixed monthly investment and the investor's tolerance for risk, how should each
instalment be split between Nifty~50 equities, gold and a liquid fund, using only the data that
would have been available at the time?}
\end{quote}

The answer must be usable by an ordinary investor every month, so the study ends not with a table
of back-test results but with a tool that gives this month's instruction.

\subsection{Objectives}

\begin{enumerate}
  \item Build a clean, verified table of monthly rupee returns for the three assets from
        Feb~2000 to the latest available month (Section~\ref{ch:data}).
  \item Build a monthly SIP simulation engine with realistic instalments, rebalancing and costs,
        in which every decision uses only past data (Section~\ref{ch:method}).
  \item Compare simple fixed mixes with optimiser-based strategies on 2010--2019, and test all of
        them, unchanged, on 2020--2025 (Section~\ref{ch:results}).
  \item Define one score that balances return against risk for a given investor, and use it to
        choose between approaches (Sections~\ref{sec:utility} and \ref{sec:res-utility}).
  \item Turn the final rule into a website that any investor can use each month
        (Section~\ref{ch:website}).
\end{enumerate}

